\documentclass[11pt]{amsart}
\usepackage{amsmath,amssymb,amsthm}
\usepackage[margin=1.15in]{geometry}
\usepackage{booktabs}
\usepackage[colorlinks=true,linkcolor=blue,citecolor=blue,urlcolor=blue]{hyperref}

\theoremstyle{plain}
\newtheorem{theorem}{Theorem}[section]
\newtheorem{proposition}[theorem]{Proposition}
\newtheorem{lemma}[theorem]{Lemma}
\newtheorem{corollary}[theorem]{Corollary}
\theoremstyle{definition}
\newtheorem{example}[theorem]{Example}
\newtheorem{problem}[theorem]{Problem}
\theoremstyle{remark}
\newtheorem{remark}[theorem]{Remark}

\newcommand{\C}{\mathbb{C}}
\newcommand{\R}{\mathbb{R}}
\newcommand{\Z}{\mathbb{Z}}
\newcommand{\ii}{\mathrm{i}}
\newcommand{\abs}[1]{\lvert #1\rvert}
\newcommand{\set}[1]{\{#1\}}
\DeclareMathOperator{\ord}{ord}

\title[Sharpness of the Enestr\"om--Kakeya annulus]{When is the Enestr\"om--Kakeya annulus sharp?\\[2pt]
Exact equality cases and a defect identity}
\author{wangyue}
\date{\today}

\begin{document}

\begin{abstract}
Let $p(z)=\sum_{k=0}^{n}a_kz^k$ with $a_k>0$, and put $A=\max_k a_k/a_{k+1}$ and
$a=\min_k a_k/a_{k+1}$. The Enestr\"om--Kakeya theorem places all zeros of $p$ in the
annulus $a\le\abs{z}\le A$. We determine exactly when this annulus is \emph{sharp}, that
is, when a zero actually lies on a boundary circle. In the standard rescaling
$b_k=a_kA^{k}$ that makes the coefficient sequence nondecreasing, a zero on the unit circle
exists if and only if $b$ is constant on consecutive blocks of some common length
$d\ge2$ dividing $n+1$; equivalently, if and only if
$q(u)=(1+u+\cdots+u^{d-1})G(u^{d})$ for such a $d$. Three consequences are that every
boundary zero is a root of unity of order dividing $n+1$, that for $n+1$ prime the outer
circle is attained only by geometric polynomials, and that a simple arithmetic condition on
the set of indices where the maximal ratio is \emph{not} attained rules out boundary zeros.
The proof is a rigidity argument for equality in the triangle inequality. It is complemented
by an exact defect identity and by a worked perturbation showing that the extremal rate in
the family $1+z+\cdots+z^{n-1}+\gamma z^{n}$ is
$2\sin^{2}(\pi/(n+1))(\gamma-1)/(n+1)=\Theta((\gamma-1)/n^{3})$. All statements are checked
computationally, over $6182$ coefficient sequences with no counterexample.
\end{abstract}

\maketitle

\section{Introduction}

\subsection{The classical statement}

Let $p(z)=\sum_{k=0}^{n}a_kz^k$ with $a_k>0$ for all $k$, and put
\begin{equation}\label{eq:ratios}
r_k:=\frac{a_k}{a_{k+1}}\quad(0\le k\le n-1),\qquad
a:=\min_{0\le k\le n-1}r_k,\qquad A:=\max_{0\le k\le n-1}r_k .
\end{equation}
The Enestr\"om--Kakeya theorem states that every zero $z$ of $p$ satisfies
\begin{equation}\label{eq:EK}
a\le\abs{z}\le A .
\end{equation}
The result is due to Enestr\"om \cite{Enestrom} and, independently and in the form above, to
Kakeya \cite{Kakeya}; see \cite[\S27]{Marden} and \cite[Chapter~4]{RS2002} for the classical
treatment and for refinements under additional hypotheses on the coefficients.

Two features of \eqref{eq:EK} are immediate. First, the annulus is optimal as stated: for
$p(z)=1+z+\cdots+z^{n}$ all ratios equal $1$ and all zeros lie on the unit circle. Second,
no annulus determined by the extreme ratios alone can be smaller, since the ratios of a
polynomial may be prescribed essentially freely.

The question studied here is different, and appears not to have been isolated in the
literature: \emph{for a given $p$, does a zero actually lie on the circle $\abs{z}=A$, or on
$\abs{z}=a$?} That is, when is \eqref{eq:EK} attained rather than merely approached? The
known refinements of \eqref{eq:EK} shrink the annulus under further hypotheses on the
coefficients; they do not characterise the polynomials for which the original annulus is
exact. Strictness has been studied for particular parametric families; the general
characterisation below covers such families as special cases.

\subsection{Main results}

The key step is the passage to the \emph{extremal normalisation}. Set
\begin{equation}\label{eq:normalisation}
b_k:=a_kA^{k}\quad(0\le k\le n),\qquad q(u):=\sum_{k=0}^{n}b_ku^{k},\qquad p(z)=q(Az).
\end{equation}
Since $b_k/b_{k+1}=r_k/A\le1$, the sequence $b$ is nondecreasing and positive, and the upper
bound in \eqref{eq:EK} becomes the statement that the zeros of $q$ lie in the closed unit
disc. Write
\begin{equation}\label{eq:Delta}
\Delta_k:=b_k-b_{k-1}\ge0\quad(1\le k\le n),\qquad
I:=\set{\,k\in\set{1,\dots,n}:\Delta_k>0\,}
\end{equation}
for the \emph{jump set} of $b$, that is, the set of indices at which the maximal ratio $A$
is not attained.

\begin{theorem}[Rigidity]\label{thm:rigidity}
Let $0<b_0\le b_1\le\cdots\le b_n$ and $q(u)=\sum_{k=0}^{n}b_ku^{k}$. The following are
equivalent.
\begin{enumerate}
\item[(i)] $q$ has a zero on the unit circle $\set{\abs{u}=1}$;
\item[(ii)] there is an integer $d\ge2$ with $d\mid(n+1)$ such that
$b_{jd}=b_{jd+1}=\cdots=b_{jd+d-1}$ for every $j=0,1,\dots,(n+1)/d-1$;
\item[(iii)] there are $d\ge2$ with $d\mid(n+1)$, an integer $m:=(n+1)/d$, and numbers
$0<c_0\le c_1\le\cdots\le c_{m-1}$ such that
\begin{equation}\label{eq:factorisation}
q(u)=\bigl(1+u+\cdots+u^{d-1}\bigr)\sum_{j=0}^{m-1}c_ju^{jd}.
\end{equation}
\end{enumerate}
\end{theorem}

Thus a boundary zero is never an accident of the coefficients: it forces a block-periodic
structure, and it always arises from the cyclotomic factor $1+u+\cdots+u^{d-1}$.

Theorem \ref{thm:rigidity} yields, in \S\ref{sec:corollaries}, that every zero on
$\abs{u}=1$ is a root of unity of order dividing $n+1$ (Corollary \ref{cor:rootofunity}),
that for $n+1$ prime the outer circle is attained only by geometric polynomials
(Corollary \ref{cor:prime}), and that the existence of boundary zeros is decided by the
arithmetic of the jump set (Corollary \ref{cor:gcd}).

Section \ref{sec:quant} gives the quantitative counterpart. There is an exact
\emph{defect identity} (Proposition \ref{prop:defect}), whose degenerate case is precisely
the rigidity statement, and which controls the deficit of $\abs{z}$ from $A$ by the
misalignment of the arguments $u^{k}$. A worked perturbation
(Proposition \ref{prop:gamma}) shows that in the family $1+z+\cdots+z^{n-1}+\gamma z^{n}$
the extremal circle is approached at the rate
$2\sin^{2}(\pi/(n+1))(\gamma-1)/(n+1)=\Theta((\gamma-1)/n^{3})$: the Enestr\"om--Kakeya
bound is sharp, but its sharpness is fragile in high degree.

Section \ref{sec:numerics} reports an exhaustive verification of
Theorem \ref{thm:rigidity} over $6182$ coefficient sequences, with no counterexample.
Section \ref{sec:open} lists the open questions.

\section{The sharp theorem and a short proof}\label{sec:EK}

For completeness, and because the proof of Theorem \ref{thm:rigidity} builds on it, we
recall a short proof of \eqref{eq:EK}. The only ingredient is the following lemma, which is
the case $A=1$ of the theorem.

\begin{lemma}[Monotone lemma]\label{lem:monotone}
Let $0<b_0\le b_1\le\cdots\le b_n$ and $q(u)=\sum_{k=0}^{n}b_ku^{k}$. Then every zero $u$
of $q$ satisfies $\abs{u}\le1$.
\end{lemma}

\begin{proof}
Expanding $(1-u)q(u)$ and collecting the coefficients of $u^{k}$ gives, for every $u$,
\[
(1-u)q(u)=b_0+\sum_{k=1}^{n}\Delta_ku^{k}-b_nu^{n+1},
\qquad \Delta_k=b_k-b_{k-1}\ge0 .
\]
Suppose that $q(u)=0$ and $R:=\abs{u}>1$. Then
$b_nu^{n+1}=b_0+\sum_{k=1}^{n}\Delta_ku^{k}$, hence
\[
b_nR^{n+1}\le b_0+\sum_{k=1}^{n}\Delta_kR^{k}
\le b_0+R^{n}\sum_{k=1}^{n}\Delta_k
=b_0+R^{n}(b_n-b_0)
=b_nR^{n}-b_0\bigl(R^{n}-1\bigr)
<b_nR^{n},
\]
because $b_0>0$ and $R>1$. Dividing by $b_nR^{n}>0$ gives $R<1$, a contradiction.
\end{proof}

\begin{theorem}[Enestr\"om--Kakeya]\label{thm:EK}
Every zero $z$ of $p$ satisfies $a\le\abs{z}\le A$.
\end{theorem}

\begin{proof}
\emph{Upper bound.} With $b_k=a_kA^{k}$ as in \eqref{eq:normalisation} we have $p(z)=q(Az)$,
where $q$ has nondecreasing positive coefficients. By Lemma \ref{lem:monotone} every zero
$u$ of $q$ satisfies $\abs{u}\le1$, so every zero $z=A^{-1}u$ of $p$ satisfies
$\abs{z}\le A$.

\emph{Lower bound.} Put $c_k:=a_ka^{k}$, so that $c_k/c_{k+1}=r_k/a\ge1$ and hence
$c_0\ge c_1\ge\cdots\ge c_n>0$. Then $p(z)=f(z/a)$ with $f(w)=\sum_{k=0}^{n}c_kw^{k}$. Let
$f^{*}(w):=w^{n}f(1/w)=\sum_{j=0}^{n}c_{n-j}w^{j}$, whose coefficients
$c_n\le c_{n-1}\le\cdots\le c_0$ are positive and nondecreasing; by Lemma \ref{lem:monotone}
every zero $w'$ of $f^{*}$ satisfies $\abs{w'}\le1$. A non-zero $w$ is a zero of $f$ if and
only if $1/w$ is a zero of $f^{*}$, so $\abs{w}\ge1$ for every zero $w$ of $f$, and
$\abs{z}=a\abs{w}\ge a$.
\end{proof}

\begin{remark}\label{rem:sharp}
The annulus \eqref{eq:EK} cannot be improved in general, as $p(z)=1+z+\cdots+z^{n}$ shows,
where $a=A=1$ and all zeros lie on $\abs{z}=1$. The content of \S\ref{sec:rigidity} is
that, for a given $p$ with $a<A$, attainment on a boundary circle is an extremely
restrictive condition.
\end{remark}

\section{Exact equality cases}\label{sec:rigidity}

We keep the notation \eqref{eq:normalisation}--\eqref{eq:Delta}: the sequence $b$ is
positive and nondecreasing, $q(u)=\sum b_ku^{k}$, and $I$ is the jump set.

\begin{lemma}[Decomposition identity]\label{lem:key}
For every $u\in\C$,
\begin{equation}\label{eq:key}
(1-u)q(u)=b_0+\sum_{k=1}^{n}\Delta_ku^{k}-b_nu^{n+1}.
\end{equation}
In particular, if $q(u)=0$ then
\begin{equation}\label{eq:star}
b_0+\sum_{k\in I}\Delta_ku^{k}=b_nu^{n+1}.
\end{equation}
\end{lemma}

\begin{proof}
Expand $(1-u)\sum_{k=0}^{n}b_ku^{k}=\sum_{k=0}^{n}b_ku^{k}-\sum_{k=0}^{n}b_ku^{k+1}$ and
collect the coefficient of $u^{k}$ for $1\le k\le n$, which is $b_k-b_{k-1}=\Delta_k$; the
constant term is $b_0$ and the term of degree $n+1$ is $-b_n$. Terms with $\Delta_k=0$ may
be dropped, which gives \eqref{eq:star}.
\end{proof}

\subsection{Proof of the rigidity theorem}

\begin{proof}[Proof of Theorem \ref{thm:rigidity}]
\emph{(i)$\Rightarrow$(ii).} Let $q(u)=0$ with $\abs{u}=1$. Since $q(1)=\sum_kb_k>0$ we
have $u\ne1$. By \eqref{eq:star},
\begin{equation}\label{eq:align0}
b_0+\sum_{k\in I}\Delta_ku^{k}=b_nu^{n+1}.
\end{equation}
Taking absolute values and using $\abs{u}=1$ together with $\sum_{k\in I}\Delta_k=b_n-b_0$,
\[
b_n=\Bigl|b_0+\sum_{k\in I}\Delta_ku^{k}\Bigr|
\le b_0+\sum_{k\in I}\Delta_k=b_n .
\]
Hence equality holds in the triangle inequality. All summands $b_0$ and $\Delta_ku^{k}$
($k\in I$) are non-zero, so they share a common argument; the summand $b_0>0$ has argument
$0$, hence
\begin{equation}\label{eq:align}
u^{k}=1\qquad\text{for every }k\in I .
\end{equation}
Substituting \eqref{eq:align} into \eqref{eq:align0} gives
$b_0+\sum_{k\in I}\Delta_k=b_nu^{n+1}$, that is $b_n=b_nu^{n+1}$, so
\begin{equation}\label{eq:rootofunity}
u^{n+1}=1 .
\end{equation}
Set $d:=\ord(u)\ge2$. By \eqref{eq:align} and \eqref{eq:rootofunity} we have $d\mid k$ for
all $k\in I$, and $d\mid(n+1)$. Since $I$ is exactly the set of indices at which $b$
strictly increases, $b$ is constant between consecutive elements of $I$; as all elements of
$I$ are multiples of $d$, the sequence $b$ is constant on every block
$\set{jd,\dots,jd+d-1}$. Because $d\mid(n+1)$, all blocks are complete, which is (ii).

\emph{(ii)$\Rightarrow$(iii).} If $b$ is constant on blocks of length $d$ and
$d\mid(n+1)$, put $c_j:=b_{jd}$ for $0\le j\le m-1$, $m=(n+1)/d$. Then
$0<c_0\le\cdots\le c_{m-1}$ and
\[
q(u)=\sum_{j=0}^{m-1}c_j\sum_{r=0}^{d-1}u^{jd+r}
=\Bigl(\sum_{r=0}^{d-1}u^{r}\Bigr)\sum_{j=0}^{m-1}c_ju^{jd},
\]
which is \eqref{eq:factorisation}.

\emph{(iii)$\Rightarrow$(i).} Let $u\ne1$ be a $d$-th root of unity. Then $\abs{u}=1$,
$1+u+\cdots+u^{d-1}=0$ and $u^{d}=1$, so $q(u)=0\cdot\sum_jc_j=0$.
\end{proof}

\begin{remark}
Only the implication (i)$\Rightarrow$(ii) uses monotonicity, and it uses it exactly twice:
to make the coefficients $\Delta_k$ non-negative in \eqref{eq:star}, and to convert the
alignment \eqref{eq:align} into the block structure of $b$. For arbitrary positive
coefficients the conclusion fails, because then the triangle inequality in the proof is
generically strict.
\end{remark}

\subsection{Corollaries}\label{sec:corollaries}

\begin{corollary}[Boundary zeros are roots of unity]\label{cor:rootofunity}
If $q(u)=0$ and $\abs{u}=1$, then $u^{n+1}=1$. Moreover $q$ vanishes at every non-trivial
$d$-th root of unity, for each $d\ge2$ dividing $n+1$ that occurs in
Theorem \ref{thm:rigidity}(ii).
\end{corollary}

\begin{proof}
The first assertion is \eqref{eq:rootofunity}; the second is the implication
(ii)$\Rightarrow$(i) of Theorem \ref{thm:rigidity}.
\end{proof}

\begin{corollary}[Prime $n+1$]\label{cor:prime}
Suppose that $n+1$ is prime. Then the circle $\abs{z}=A$ contains a zero of $p$ if and only
if all ratios $r_k$ coincide, that is $a=A$ and $p(z)=c\sum_{k=0}^{n}(z/A)^{k}$ for some
$c>0$. In that case \eqref{eq:EK} degenerates to a single circle.
\end{corollary}

\begin{proof}
If a zero with $\abs{z}=A$ exists, Theorem \ref{thm:rigidity} provides $d\ge2$ with
$d\mid(n+1)$; primality forces $d=n+1$, so $b$ is constant on the single block
$\set{0,\dots,n}$, i.e.\ $b_0=\cdots=b_n$. Then $b_k=b_{k+1}$ gives $r_k=A$ for all $k$,
hence $a=A$ and $a_k=a_0A^{-k}$. Conversely, if $a_k=a_0A^{-k}$ then
$p(z)=a_0\sum_k(z/A)^{k}$, whose zeros satisfy $(z/A)^{n+1}=1$ with $z\ne A$, hence
$\abs{z}=A$.
\end{proof}

\begin{corollary}[Arithmetic criterion]\label{cor:gcd}
Let $I$ be the jump set \eqref{eq:Delta} and let $g$ be the greatest common divisor of the
elements of $I$ (with $g:=0$ if $I=\emptyset$).
\begin{enumerate}
\item[(a)] If $I=\emptyset$, equivalently $b_0=\cdots=b_n$ and $a=A$, then $q$ vanishes at
all $(n+1)$-th roots of unity other than $1$.
\item[(b)] If $I\ne\emptyset$, then $q$ has a zero on the unit circle if and only if
$\gcd(g,n+1)\ge2$.
\end{enumerate}
In particular, if the maximal ratio $A$ is attained at a single index $k_0$ and
$\gcd(k_0,n+1)=1$, then no zero of $p$ lies on $\abs{z}=A$.
\end{corollary}

\begin{proof}
Part (a) is the case $d=n+1$ of Theorem \ref{thm:rigidity}, in which the single block
covers all coefficients. For (b), Theorem \ref{thm:rigidity} says that a boundary zero
exists if and only if $b$ is constant on blocks of some length $d\ge2$ with $d\mid(n+1)$;
this holds if and only if $d$ divides every element of $I$, that is $d\mid g$. Such a
$d\ge2$ exists if and only if $\gcd(g,n+1)\ge2$. The last assertion is the case
$I=\set{k_0}$.
\end{proof}

\begin{example}\label{ex:examples}
\leavevmode
\begin{enumerate}
\item[(a)] $q(u)=1+u+2u^{2}+2u^{3}$, so $n+1=4$, $I=\set2$ and $\gcd(2,4)=2$. Boundary
zeros exist: indeed $q(u)=(1+u)(1+2u^{2})$, and $u=-1$ has modulus $1$. The remaining zeros
$\pm\ii/\sqrt2$ have modulus $1/\sqrt2<1$, so the outer circle is attained and the inner one
is not.
\item[(b)] $q(u)=1+u+u^{2}+\gamma u^{3}$ with $\gamma>1$, so $n+1=4$ and $I=\set3$. Since
$\gcd(3,4)=1$, there is no zero of modulus $1$: all zeros lie strictly inside the unit
disc. Note that $b=(1,1,1,\gamma)$ is \emph{almost} block-periodic with $d=2$, the only
failure being in the last block, and that this already destroys attainment.
\item[(c)] $q(u)=1+u+u^{2}+u^{3}+2u^{4}$, so $n+1=5$ is prime. There is no zero on the unit
circle, in agreement with Corollary \ref{cor:prime}.
\item[(d)] Take $b=(1,1,2,3,3,3)$, that is
$q(u)=1+u+2u^{2}+3u^{3}+3u^{4}+3u^{5}$, with $n+1=6$. Here $I=\set{2,3}$ and
$\gcd(\gcd(2,3),6)=1$, so there is no zero of modulus $1$; numerically the maximal modulus
is $0.878335\ldots$. Note that $b$ is a small perturbation of $(1,1,1,3,3,3)$, which
\emph{is} block-periodic with $d=3$ and does attain the bound. A single misplaced
coefficient therefore destroys attainment, in agreement with the rigidity of
Theorem \ref{thm:rigidity}.
\end{enumerate}
\end{example}

\begin{remark}[The inner circle]\label{rem:dual}
Attainment of the lower bound is governed by the reversed sequence. With $c_k=a_ka^{k}$ and
$f(w)=\sum_kc_kw^{k}$ as in the proof of Theorem \ref{thm:EK}, a zero $z$ of $p$ satisfies
$\abs{z}=a$ if and only if the zero $w=z/a$ of $f$ satisfies $\abs{w}=1$, that is, if and
only if $f^{*}(w)=\sum_jc_{n-j}w^{j}$ satisfies the hypothesis of
Theorem \ref{thm:rigidity} with $c_{n-j}$ in place of $b_j$. Reversing a block structure of
length $d$ with $d\mid(n+1)$ produces a block structure of the same length, so the two
conditions have the same shape, but they are imposed on the two different rescalings
$a_kA^{k}$ and $a_ka^{k}$; a sequence may well attain one bound and not the other, as in
Example \ref{ex:examples}(a).
\end{remark}

\section{The quantitative side}\label{sec:quant}

Theorem \ref{thm:rigidity} describes the exact equality case. The next two results quantify
what happens nearby: an identity that measures the deficit, and a computation that shows how
fast the extremal circle is approached under a natural perturbation.

\subsection{A defect identity}

\begin{proposition}[Defect identity]\label{prop:defect}
Let $0<b_0\le\cdots\le b_n$ and let $u$ be a zero of $q$. Then
\begin{equation}\label{eq:defect}
b_n\bigl(1-\operatorname{Re}u^{n+1}\bigr)
=\sum_{k=1}^{n}\Delta_k\bigl(1-\operatorname{Re}u^{k}\bigr).
\end{equation}
Consequently, with $R=\abs{u}$ and $u=Re^{\ii\vartheta}$,
\begin{equation}\label{eq:defect2}
\sum_{k\in I}\Delta_k\bigl(1-R^{k}\cos(k\vartheta)\bigr)\ \ge\ b_n\bigl(1-R^{n+1}\bigr).
\end{equation}
\end{proposition}

\begin{proof}
Taking real parts in \eqref{eq:star} and using $\sum_{k\in I}\Delta_k=b_n-b_0$ gives
\[
b_n\operatorname{Re}u^{n+1}
=b_0+\sum_{k\in I}\Delta_k\operatorname{Re}u^{k}
=b_n-\sum_{k\in I}\Delta_k+\sum_{k\in I}\Delta_k\operatorname{Re}u^{k},
\]
which is \eqref{eq:defect}. For \eqref{eq:defect2}, use
$\operatorname{Re}u^{n+1}\le R^{n+1}$ in the left-hand side of \eqref{eq:defect} and
$\operatorname{Re}u^{k}=R^{k}\cos(k\vartheta)$ in the right-hand side.
\end{proof}

\begin{remark}\label{rem:defectrigid}
At $R=1$ the identity \eqref{eq:defect} reads
$b_n\bigl(1-\cos((n+1)\vartheta)\bigr)=\sum_{k\in I}\Delta_k\bigl(1-\cos(k\vartheta)\bigr)$,
with both sides non-negative. This is a ``linearised'' form of the rigidity argument, but it
is strictly weaker: the real part alone does not force the individual terms to vanish, which
is why the alignment \eqref{eq:align} was deduced from the modulus rather than from
\eqref{eq:defect}.
\end{remark}

\subsection{How fast is the extremal circle approached?}

Fix $n\ge1$ and consider the family
\begin{equation}\label{eq:gamma}
p_\gamma(z)=1+z+\cdots+z^{n-1}+\gamma z^{n},\qquad \gamma>1 .
\end{equation}
Its ratios are $r_0=\cdots=r_{n-2}=1$ and $r_{n-1}=1/\gamma$, so $A=1$, $a=1/\gamma$ and the
jump set is $I=\set n$. By Corollary \ref{cor:gcd}(b) there is no zero of modulus $1$,
since $\gcd(n,n+1)=1$. It is therefore natural to ask how far the zeros are from the unit
circle.

\begin{proposition}[Sharpness rate]\label{prop:gamma}
Let $\omega=e^{2\pi\ii k/(n+1)}$ with $1\le k\le n$, and let $z_k(\gamma)$ denote the zero
of $p_\gamma$ near $\omega$. Then $z_k$ is analytic near $\gamma=1$ and
\begin{equation}\label{eq:rate}
z_k(\gamma)=\omega+\frac{1-\omega}{n+1}(\gamma-1)+O\bigl((\gamma-1)^{2}\bigr),
\end{equation}
\begin{equation}\label{eq:rate2}
\abs{z_k(\gamma)}
=1-\frac{2\sin^{2}\bigl(\pi k/(n+1)\bigr)}{n+1}(\gamma-1)+O\bigl((\gamma-1)^{2}\bigr).
\end{equation}
In particular
\begin{align}
\max_{1\le k\le n}\abs{z_k(\gamma)}
&=1-\frac{2\sin^{2}\bigl(\pi/(n+1)\bigr)}{n+1}(\gamma-1)+O\bigl((\gamma-1)^{2}\bigr)\label{eq:rate3}\\
&=1-\frac{2\pi^{2}}{(n+1)^{3}}(\gamma-1)+O\bigl((\gamma-1)^{2}\bigr)
+O\!\left(\frac{\gamma-1}{(n+1)^{5}}\right).\label{eq:rate4}
\end{align}
\end{proposition}

\begin{proof}
Write $p_\gamma(z)=\dfrac{1-z^{n+1}}{1-z}+(\gamma-1)z^{n}$ and set $\varepsilon=\gamma-1$.
At $\varepsilon=0$ the zeros are the $(n+1)$-th roots of unity $\omega\ne1$, all simple, so
by the implicit function theorem each of them continues analytically to a zero
$z_k(\varepsilon)$ of $p_\gamma$. Differentiating $p_\gamma(z_k(\varepsilon))=0$ at
$\varepsilon=0$ gives $z_k'(0)=-\omega^{n}/p_1'(\omega)$. Now
$p_1'(z)=\sum_{k=1}^{n}kz^{k-1}=\dfrac{1-(n+1)z^{n}+nz^{n+1}}{(1-z)^{2}}$, so at
$z=\omega$, using $\omega^{n+1}=1$ and $\omega\ne1$,
\[
p_1'(\omega)=\frac{1-(n+1)\omega^{-1}+n}{(1-\omega)^{2}}
=\frac{(n+1)\bigl(1-\omega^{-1}\bigr)}{(1-\omega)^{2}}
=-\frac{(n+1)\omega^{-1}}{1-\omega},
\]
and therefore
\[
z_k'(0)=-\frac{\omega^{n}}{p_1'(\omega)}
=\omega^{n}\cdot\frac{1-\omega}{(n+1)\omega^{-1}}
=\frac{\omega^{n+1}(1-\omega)}{n+1}
=\frac{1-\omega}{n+1},
\]
which is \eqref{eq:rate}. Then
\[
\abs{z_k(\varepsilon)}^{2}
=1+\frac{2\varepsilon}{n+1}\operatorname{Re}\bigl(\overline{\omega}(1-\omega)\bigr)
+O(\varepsilon^{2})
=1-\frac{4\varepsilon}{n+1}\sin^{2}\!\Bigl(\frac{\pi k}{n+1}\Bigr)+O(\varepsilon^{2}),
\]
since $\operatorname{Re}\bigl(\overline{\omega}(1-\omega)\bigr)=\cos\vartheta-1
=-2\sin^{2}(\vartheta/2)$ with $\vartheta=2\pi k/(n+1)$. Taking square roots gives
\eqref{eq:rate2}. The function $k\mapsto\sin^{2}(\pi k/(n+1))$ is maximal over
$1\le k\le n$ at $k=1$, which gives the first equality in \eqref{eq:rate3}; the second
follows from $\sin t=t+O(t^{3})$.
\end{proof}

\begin{remark}
Proposition \ref{prop:gamma} makes the fragility of \eqref{eq:EK} quantitative: a relative
perturbation of size $\varepsilon$ moves the extremal zeros off the circle by only
$\Theta(\varepsilon/n^{3})$, so forcing a zero to lie a distance $\delta$ inside the disc
requires $\gamma-1\asymp n^{3}\delta$. The exponent $3$ reflects the double degeneracy at
$u=1$, where both the block structure and the $(n+1)$-th root of unity condition collapse.
\end{remark}

\begin{table}[ht]
\centering
\caption{Proposition \ref{prop:gamma} against direct computation of the maximal root modulus
of \eqref{eq:gamma}.}
\label{tab:gamma}
\begin{tabular}{cccc}
\toprule
$n$ & $\gamma$ & $\max_k\abs{z_k(\gamma)}$ & first-order prediction \\
\midrule
$4$  & $1.001$ & $0.999861815004$ & $0.999861803399$ \\
$4$  & $1.1$   & $0.986330205396$ & $0.986180339887$ \\
$9$  & $1.001$ & $0.999980891439$ & $0.999980901699$ \\
$9$  & $1.1$   & $0.997994875902$ & $0.998090169944$ \\
$19$ & $1.001$ & $0.999997550857$ & $0.999997552826$ \\
$19$ & $1.1$   & $0.999736088306$ & $0.999755282581$ \\
\bottomrule
\end{tabular}
\end{table}

\section{Computational verification}\label{sec:numerics}

All statements of \S\ref{sec:rigidity} and \S\ref{sec:quant} were checked with the scripts
in the companion repository (\texttt{verify/}); the entry points are \texttt{ek\_core.py},
\texttt{ek\_more.py} and \texttt{ek\_dual.py}.

\paragraph{Method.} For an integer sequence $b$, the question whether $q$ has a zero on
$\abs{u}=1$ is decided as follows. Since $q$ has real coefficients, a zero $u$ with
$\abs{u}=1$ satisfies $q(u)=u^{n}\overline{q(u)}=q^{*}(u)$, where $q^{*}(u)=u^{n}q(1/u)$ is
the reversed polynomial. Hence a boundary zero forces the gcd $\gcd(q,q^{*})$ to be
non-constant, and this is decided \emph{exactly} in rational arithmetic: when
$\gcd(q,q^{*})=1$ there is no boundary zero and no floating-point computation enters the
decision. When the gcd is non-constant, its zeros are computed numerically and their moduli
inspected; in the search below this second step always produced a root of unity, in
agreement with Corollary \ref{cor:rootofunity}.

\paragraph{Exhaustive search.} We enumerated all nondecreasing sequences
$1\le b_0\le\cdots\le b_n\le5$ with $2\le n+1\le12$, that is $6182$ sequences, and compared
the exact decision with the block-structure criterion of Theorem \ref{thm:rigidity}(ii).
There were $591$ sequences with a boundary zero, all of them block-structured, and no
counterexample in either direction. In the same sweep, all $1111$ zeros found on the unit
circle satisfied $u^{n+1}=1$, confirming Corollary \ref{cor:rootofunity}.

\paragraph{The two-sided question.} Among the $6116$ sequences of the sweep with $a<A$, that
is with a non-degenerate annulus, none attained both boundary circles. This is consistent
with, but does not prove, the expectation discussed in Problem \ref{prob:both}.

\paragraph{Defect identity.} The identity \eqref{eq:defect} was verified numerically on
$2000$ random polynomials with positive rational coefficients of degree at most $8$; the
maximal relative error was $1.2\cdot10^{-14}$, at the level of round-off. Table
\ref{tab:gamma} confirms Proposition \ref{prop:gamma}.

\section{Open questions}\label{sec:open}

\begin{problem}[Quantitative rigidity]\label{prob:quant}
Find an explicit functional $\eta$ of the coefficient sequence, vanishing exactly on the
block-structured sequences of Theorem \ref{thm:rigidity}(ii), such that every zero $u$ of
$q$ satisfies $\abs{u}\le1-\eta(b)$ for a non-degenerate annulus. Proposition
\ref{prop:gamma} shows that the true deficit can be as small as $\Theta(\varepsilon/n^{3})$
for a perturbation of relative size $\varepsilon$, which suggests that $\eta$ must be cubic
in the distance to block structure and that a linear stability estimate cannot hold.
\end{problem}

\begin{problem}[Simultaneous attainment]\label{prob:both}
Suppose that $a<A$ and that both circles $\abs{z}=a$ and $\abs{z}=A$ contain a zero of $p$.
Must then $a=A$? The search of \S\ref{sec:numerics} found no counterexample. The difficulty
is that the two conditions of Theorem \ref{thm:rigidity} are imposed on the two different
rescalings $a_kA^{k}$ and $a_ka^{k}$ of the same coefficient sequence.
\end{problem}

\begin{problem}[Other settings]\label{prob:other}
The rigidity argument uses only the ordered field structure and the triangle inequality, so
analogues should hold for the quaternionic, octonionic and matrix versions of the
Enestr\"om--Kakeya theorem that have been studied in the literature. A discrete analogue for
the polynomial-method bounds on Kakeya sets in finite fields also seems plausible. We did
not pursue either direction.
\end{problem}

\begin{thebibliography}{9}

\bibitem[E]{Enestrom}
G.~Enestr\"om,
\emph{Remarque sur un th\'eor\`eme relatif aux racines de l'\'equation
$a_nx^{n}+\cdots+a_0=0$ o\`u tous les coefficients sont r\'eels et positifs},
T\^ohoku Math.\ J.\ \textbf{18} (1920), 34--36.

\bibitem[K]{Kakeya}
S.~Kakeya,
\emph{On the limits of the roots of an algebraic equation with positive coefficients},
T\^ohoku Math.\ J.\ \textbf{2} (1912), 140--142.

\bibitem[M]{Marden}
M.~Marden,
\emph{Geometry of Polynomials}, second ed.,
Mathematical Surveys and Monographs, no.~3, American Mathematical Society, Providence, RI,
1966.

\bibitem[RS]{RS2002}
Q.~I. Rahman and G.~Schmeisser,
\emph{Analytic Theory of Polynomials},
London Mathematical Society Monographs (New Series), no.~26, Oxford University Press,
Oxford, 2002.

\end{thebibliography}

\end{document}
